\chapter{Manual técnico}

Este manual recolle a información necesaria para instalar, executar, probar e modificar
Telaria. Diríxese a persoas que vaian continuar o desenvolvemento do sistema.

\section{Requisitos do entorno}

\begin{itemize}
  \item \textbf{JDK 25} (para o backend; a construción realízase co \textit{wrapper} de Gradle).
  \item \textbf{Node.js} con \texttt{npm} (para o frontend e a xeración de tipos).
  \item \textbf{Podman} con soporte de \texttt{podman compose} (para os servizos de apoio).
  \item \textbf{Git}.
\end{itemize}

\section{Estrutura do repositorio}

\begin{itemize}
  \item \texttt{backend/} — API REST en Spring Boot (Gradle).
  \item \texttt{frontend/} — aplicación móbil en React Native + Expo.
  \item \texttt{memoria/} — esta documentación (LaTeX).
\end{itemize}

\section{Servizos de apoio (Podman)}

Os servizos auxiliares (PostgreSQL, MinIO e Ollama) defínense en
\texttt{backend/compose.yaml}. O xeito máis sinxelo de poñelos en marcha, descargar o
modelo de linguaxe e arrincar o backend é o \textit{script} de inicio:

\begin{lstlisting}[language=bash]
cd backend
./start.sh
\end{lstlisting}

\noindent Alternativamente, pódense levantar só os contedores:

\begin{lstlisting}[language=bash]
cd backend
podman compose up -d        # engadir --profile gpu se hai GPU dispoñible
\end{lstlisting}

\noindent Portos por defecto: PostgreSQL \texttt{5432}, MinIO \texttt{9000}/\texttt{9001}
(API/consola), Ollama \texttt{11434}.

\section{Execución do backend}

\begin{lstlisting}[language=bash]
cd backend
./gradlew bootRun
\end{lstlisting}

\noindent O servidor escoita no porto \texttt{8082}. A configuración (base de datos,
\textit{keystore} de sinatura JWT, MinIO e tarefas programadas) atópase en
\texttt{src/main/resources/application.yaml}. A documentación interactiva da API queda
dispoñible en \texttt{/swagger-ui.html}.

\section{Execución do frontend}

\begin{lstlisting}[language=bash]
cd frontend
npm install
npm start            # inicia o servidor de desenvolvemento de Expo
\end{lstlisting}

\noindent A dirección do backend configúrase na constante \texttt{BASE\_URL} de
\texttt{frontend/constants/constants.tsx}, que debe apuntar ao enderezo accesible do
servidor desde o dispositivo. A aplicación pode abrirse en Expo Go ou nun emulador
(\texttt{npm run android} / \texttt{npm run ios}).

\section{Contrato OpenAPI e xeración de código}

A API especifícase en \texttt{backend/src/main/resources/openapi/}, cun documento
principal (\texttt{api.yaml}) que referencia ficheiros de rutas e esquemas. A partir del
xérase código:

\begin{lstlisting}[language=bash]
# Backend: interfaces de controladores e DTO
cd backend && ./gradlew openApiGenerate

# Frontend: tipos TypeScript
cd frontend && npm run generate:types
\end{lstlisting}

\noindent Ao editar ficheiros referenciados con \texttt{\$ref} (rutas ou esquemas), pode
ser preciso forzar a rexeneración, xa que a tarefa de Gradle só observa cambios en
\texttt{api.yaml}.

\section{Execución das probas}

As probas de integración (E2E) empregan Testcontainers, que require un \textit{socket}
compatible con Docker. Con Podman \textit{rootless}, débese indicar a súa ruta:

\begin{lstlisting}[language=bash]
cd backend
export DOCKER_HOST=unix:///run/user/$(id -u)/podman/podman.sock
export TESTCONTAINERS_RYUK_DISABLED=true
./gradlew test
\end{lstlisting}

\noindent Ao rematar, o informe de cobertura JaCoCo xérase en
\texttt{backend/build/reports/jacoco/test/html/index.html}. A análise estática SpotBugs
está dispoñible mediante a tarefa correspondente de Gradle.

\section{Configuración e segredos}

A configuración do backend (credenciais da base de datos, contrasinal do \textit{keystore}
JWT, claves de acceso a MinIO e parámetros das tarefas programadas) reside en
\texttt{application.yaml}. Para un despregamento en produción recoméndase externalizar
estes valores mediante variables de entorno e non incluílos no control de versións.

\section{Integración continua}

O repositorio inclúe fluxos de traballo de GitHub Actions que executan a construción e as
probas do backend, así como unha análise de seguridade do código (CodeQL).
